\section{CTMS}

% SLIDE ID: sec06-goal
\BlockGoalFrame{\SlideID{sec06-goal}}{
  \item zrozumieć CTMS jako system zarządzania operacyjnego badaniem,
  \item odróżnić informacje operacyjne od danych klinicznych w EDC,
  \item zobaczyć, jak CTMS wspiera ośrodki, monitoring i oversight sponsora/CRO.
}

% SLIDE ID: sec06-core
\begin{frame}{CTMS jako centrum operacyjne badania}
  \SlideID{sec06-core}
  \begin{itemize}
    \item CTMS służy do planowania, śledzenia i nadzoru nad przebiegiem badania,
    \item koncentruje się na operacjach, a nie na pojedynczych danych klinicznych,
    \item łączy informacje o ośrodkach, rekrutacji, monitoringu, finansach i statusach działań,
    \item dla sponsora lub \texttt{CRO} staje się operacyjnym „single source of truth” dla portfolio badań.
  \end{itemize}
\end{frame}

% SLIDE ID: sec06-history-overview
\begin{frame}{Historia rozwoju CTMS}
  \SlideID{sec06-history-overview}
  \begin{center}
    \Large
    Arkusze i papier $\rightarrow$ lokalne bazy danych $\rightarrow$ wyspecjalizowane systemy \texttt{CTMS} $\rightarrow$ chmura i ekosystem \texttt{eClinical}
  \end{center}

  \vspace{0.45cm}

  \begin{itemize}
    \item Historia \texttt{CTMS} jest historią porządkowania coraz bardziej złożonych operacji badania klinicznego.
    \item Wraz ze wzrostem liczby ośrodków, krajów, dostawców i wymagań regulacyjnych przestały wystarczać ręczne listy, lokalne pliki i rozproszone raporty.
    \item Współczesny \texttt{CTMS} nie jest już prostym trackerem zadań, ale szkieletem operacyjnym całego badania.
  \end{itemize}
\end{frame}

% SLIDE ID: sec06-history-pre-digital
\begin{frame}{1. Era przedcyfrowa i wczesne bazy: lata 80. i 90.}
  \SlideID{sec06-history-pre-digital}
  \begin{itemize}
    \item Zanim pojawiły się wyspecjalizowane systemy, zespoły badawcze pracowały na arkuszach kalkulacyjnych, papierowych listach i lokalnych zestawieniach.
    \item Wraz ze wzrostem skali badań zaczęto budować proste bazy \texttt{Access} lub \texttt{SQL}, które dawały podstawowe przechowywanie danych operacyjnych.
    \item Problemem było jednak rozproszenie informacji: osobno plan wizyt, osobno status ośrodka, osobno płatności, osobno dokumenty i milestone'y.
    \item Taki model był silnie ręczny, podatny na duplikację wpisów i praktycznie nie dawał bieżącej widoczności sytuacji w wielu ośrodkach jednocześnie.
  \end{itemize}
\end{frame}

% SLIDE ID: sec06-history-purpose-built
\begin{frame}{2. Powstanie wyspecjalizowanych CTMS: koniec lat 90. i lata 2000}
  \SlideID{sec06-history-purpose-built}
  \begin{itemize}
    \item Wraz z harmonizacją wymagań \texttt{ICH-GCP} w 1996 r. rosła potrzeba lepszego śledzenia działań operacyjnych, nadzoru i audytowalności procesów badania.
    \item Zaczęły pojawiać się komercyjne, purpose-built systemy \texttt{CTMS}, projektowane specjalnie do planowania badania, pracy z ośrodkami, monitoringu i budżetów.
    \item Ich przewagą była standaryzacja: te same statusy, te same workflow i bardziej spójny obraz portfolio badań dla sponsora lub \texttt{CRO}.
    \item To był moment przejścia z „pliku do zarządzania badaniem” do prawdziwego systemu operacyjnego dla badań klinicznych.
  \end{itemize}
\end{frame}

% SLIDE ID: sec06-history-cloud
\begin{frame}{3. Przejście do chmury i integracji eClinical: lata 2010}
  \SlideID{sec06-history-cloud}
  \begin{itemize}
    \item W latach 2010 coraz więcej systemów \texttt{CTMS} przeszło z modelu \texttt{on-premises} do \texttt{cloud / SaaS}, co ułatwiło współpracę sponsora, \texttt{CRO} i zdalnych zespołów.
    \item Coraz wyraźniej dostrzegano problem silosów: osobny \texttt{EDC}, osobny \texttt{eTMF}, osobny system start-up, osobne raporty i ręczne reconciliacje.
    \item Z tego powodu zaczęła rosnąć rola „connected ecosystem”, czyli integracji \texttt{CTMS} z \texttt{EDC}, \texttt{eTMF}, randomizacją, safety i finansami.
    \item W praktyce oznaczało to mniej duplikacji danych, lepszą widoczność statusów i szybsze decyzje operacyjne.
  \end{itemize}
\end{frame}

% SLIDE ID: sec06-history-modern
\begin{frame}{4. Współczesny CTMS: lata 2020 do dziś}
  \SlideID{sec06-history-modern}
  \begin{itemize}
    \item Dzisiejszy \texttt{CTMS} działa jak cyfrowe rusztowanie badania: pokazuje start-up, status ośrodków, rekrutację, wizyty monitorujące, płatności i operacyjne \texttt{KPI}.
    \item Duże znaczenie mają dashboardy czasu rzeczywistego, automatyczne alerty, audit trail i gotowość do inspekcji lub audytu.
    \item W nowoczesnych platformach rośnie znaczenie centralnego oversightu sponsora, pracy z modelami zdalnymi i wspierania \texttt{risk-based monitoring}.
    \item Kierunek rozwoju nie polega już tylko na „więcej funkcji”, ale na lepszym łączeniu operacji, zgodności i informacji zarządczej.
  \end{itemize}
\end{frame}

% SLIDE ID: sec06-history-what-changed
\begin{frame}{Co zmienił CTMS w praktyce}
  \SlideID{sec06-history-what-changed}
  \begin{itemize}
    \item Zmienił widoczność: sponsor i \texttt{CRO} nie muszą już czekać na ręczne raporty z wielu źródeł.
    \item Zmienił odpowiedzialność: łatwiej wskazać właściciela zadania, otwartą akcję, opóźniony milestone lub ryzyko w ośrodku.
    \item Zmienił tempo decyzji: dashboard i aktualny status pozwalają reagować wcześniej niż w modelu opartym na mailach i arkuszach.
    \item Zmienił też relację z innymi systemami: \texttt{CTMS} staje się systemem zarządzania operacyjnego, a nie magazynem danych klinicznych.
  \end{itemize}
\end{frame}

% SLIDE ID: sec06-vs-edc
\begin{frame}{CTMS a EDC}
  \SlideID{sec06-vs-edc}
  \begin{columns}[T,onlytextwidth]
    \column{0.48\textwidth}
    \begin{beamercolorbox}[wd=\linewidth,sep=0.8em,rounded=true]{block body}
      \textbf{\textcolor{UAFMTeal}{EDC}}\par\medskip
      \begin{itemize}
        \item dane pacjenta,
        \item formularze i wizyty,
        \item walidacje i query.
      \end{itemize}
    \end{beamercolorbox}

    \column{0.48\textwidth}
    \begin{beamercolorbox}[wd=\linewidth,sep=0.8em,rounded=true]{block body}
      \textbf{\textcolor{UAFMDarkBlue}{CTMS}}\par\medskip
      \begin{itemize}
        \item ośrodki i plan badania,
        \item monitoring i działania zespołu,
        \item statusy, KPI i issue tracking.
      \end{itemize}
    \end{beamercolorbox}
  \end{columns}
\end{frame}

% SLIDE ID: sec06-sites
\begin{frame}{Zarządzanie ośrodkami i rekrutacją}
  \SlideID{sec06-sites}
  \begin{itemize}
    \item status uruchomienia ośrodka,
    \item planowana i rzeczywista rekrutacja,
    \item gotowość do pierwszego pacjenta,
    \item porównanie ośrodków pod względem aktywności i tempa pracy.
  \end{itemize}
\end{frame}

% SLIDE ID: sec06-monitoring
\begin{frame}{Jak CTMS wspiera monitoring}
  \SlideID{sec06-monitoring}
  \begin{itemize}
    \item planuje wizyty monitorujące i follow-up,
    \item porządkuje listę otwartych zadań i obserwacji,
    \item wspiera przygotowanie raportów i kolejnych działań w ośrodku.
  \end{itemize}
\end{frame}

% SLIDE ID: sec06-statuses
\begin{frame}{Statusy badania, KPI i issue tracking}
  \SlideID{sec06-statuses}
  \begin{itemize}
    \item statusy pokazują, gdzie jest badanie i każdy ośrodek,
    \item KPI wspierają oversight nad rekrutacją, monitoringiem i terminowością,
    \item issue tracking porządkuje problemy, właścicieli i terminy działań.
  \end{itemize}
\end{frame}

% SLIDE ID: sec06-pm-view
\begin{frame}{Co widzi Project Manager?}
  \SlideID{sec06-pm-view}
  \begin{itemize}
    \item ogólny status badania i ośrodków,
    \item postęp rekrutacji i kamienie milowe,
    \item opóźnienia, ryzyka i otwarte problemy,
    \item poziom realizacji działań przez zespół i CRO.
  \end{itemize}
\end{frame}

% SLIDE ID: sec06-cra-view
\begin{frame}{Co widzi CRA?}
  \SlideID{sec06-cra-view}
  \begin{itemize}
    \item listę ośrodków i plan wizyt,
    \item otwarte działania po poprzednim monitoringu,
    \item status dokumentów i gotowość ośrodka,
    \item kwestie wymagające eskalacji lub follow-up.
  \end{itemize}
\end{frame}

% SLIDE ID: sec06-oversight
\begin{frame}{Oversight sponsora / CRO}
  \SlideID{sec06-oversight}
  \begin{itemize}
    \item sponsor lub CRO widzi badanie z poziomu operacyjnego,
    \item może śledzić wykonanie działań bez wchodzenia w każdy szczegół kliniczny,
    \item CTMS wspiera kontrolę terminów, jakości pracy i eskalacji problemów,
    \item nawet przy outsourcingu odpowiedzialność za oversight nie znika, dlatego centralny widok operacyjny ma znaczenie regulacyjne i zarządcze.
  \end{itemize}
\end{frame}

% SLIDE ID: sec06-demo
\DemoFrame{Status ośrodka i wizyta monitorująca}{
  widok ośrodka, status rekrutacji, plan i wykonanie wizyty monitorującej, otwarte akcje oraz najważniejsze wskaźniki operacyjne.
}

% SLIDE ID: sec06-screens
\begin{frame}{Jakie widoki CTMS warto pokazać}
  \SlideID{sec06-screens}
  \begin{itemize}
    \item dashboard badania z kamieniami milowymi, statusem krajów i ośrodków,
    \item karta ośrodka z datami aktywacji, rekrutacją i planem wizyt,
    \item kalendarz lub lista wizyt monitorujących wraz z \texttt{follow-up actions},
    \item ekran \texttt{issue tracking} z właścicielem, terminem i statusem eskalacji,
    \item widok płatności lub budżetu, jeśli system to wspiera,
    \item raport \texttt{KPI} pokazujący opóźnienia, produktywność ośrodków i ryzyka operacyjne.
  \end{itemize}
\end{frame}

% SLIDE ID: sec06-discussion
\DiscussionFrame{Które informacje powinny trafiać do CTMS, a które nie?}{
  Gdzie kończy się obszar operacyjny \texttt{CTMS}, a zaczyna obszar danych klinicznych w \texttt{EDC}? Jakie ryzyko pojawia się wtedy, gdy zespół miesza te role albo próbuje używać jednego systemu do wszystkiego?
}

% SLIDE ID: sec06-sources
\begin{frame}{Źródła do tej sekcji}
  \SlideID{sec06-sources}
  \footnotesize
  \begin{itemize}
    \item Clinion, \textit{What is a CTMS? A Complete Guide for Researchers}, 29 grudnia 2025.
    \item PHARMASEAL, \textit{What is a CTMS - Clinical Trial Management System?}
    \item ICH, \textit{Guideline for Good Clinical Practice E6(R1)}, Step 4 dated 10 June 1996.
    \item FDA, \textit{E6(R3) Good Clinical Practice (GCP)}, final guidance, September 2025.
    \item EMA, \textit{ICH E6 Good clinical practice - Scientific guideline}, timeline dla \texttt{E6(R3)}.
  \end{itemize}
\end{frame}
